\chapter[Continuous families and exact bounds]{Continuous history families: projection, optimization, and threshold refinement}
\label{sim:continuous-methods}

A finite list of histories is often an inadequate description even of a small company. Unknown rate exposure, customer allocation, asset age and settlement timing can vary continuously. The right question is frequently whether a conclusion holds over a constrained region. This chapter develops several exact small calculations and shows where they cease to be adequate. The methods are proposals for bounded profiles, not a claim that every company reconstruction is a linear program.

\section{An affine exposure family}
Let $x$ be foreign-currency debt maturing soon and $a$ the floating share of total debt. Suppose disclosed margins imply $0\leq x\leq50$ and $0\leq a\leq1$. A simplified next-period cash requirement under currency depreciation $d$ and rate increase $r$ is
\begin{equation}
 Q(x,a)=Q_0+dx+100ra.
\end{equation}
If $d$ and $r$ are nonnegative and the two unknowns are jointly unconstrained beyond their bounds, the maximum is $Q_0+50d+100r$ and minimum is $Q_0$. These extrema occur at feasible corners. No grid is needed.

Now add a disclosure that the floating portion is entirely within the foreign-currency debt maturing soon, so $100a\leq x$. The feasible family is no longer a rectangle. The minimum remains $Q_0$, but the maximum floating share is one half, giving maximum $Q_0+50d+50r$. The relation changes the worst case even though debt total remains 100.

The condition is a stipulated example, not an inference from the margins alone. It illustrates why exact bounds must optimize over the joint constraints. Independently maximizing every marginal amount can overstate or understate the relevant result depending on signs and relationships. A constraint-aware bound carries the assumptions that make its corner feasible.

If the query includes a hedge whose receipt increases with $x$, the coefficient on $x$ can change sign. Monotonicity should be derived from the full expression, not assumed from a variable's label. A larger foreign-currency exposure can have a smaller net cash effect when a linked hedge changes at the same time.

\section{Linear programs as a bounded exact method}
When observations, bridges and query are linear in continuous unknowns, a bound can be formulated as a linear program. Let $z$ collect amounts and let $Az\leq b$, $Ez=e$ define the family. The minimum cash query is
\begin{equation}
 m=\inf\{c^{\mathsf T}z+c_0:Az\leq b, Ez=e\}.
\end{equation}
The formulation separates a mathematical target from a solver. A solver can return an optimum, infeasibility, unboundedness or an unresolved numerical status. Each requires interpretation. Unboundedness may reveal a missing economic constraint rather than a truly unlimited company exposure.

A dual certificate can support a bound when its assumptions and numerical accuracy are checked. For a simple example, constraints $x\geq0$, $y\geq0$ and $x+y\leq10$ imply $2x+y\leq20$ because $2x+y\leq2(x+y)\leq20$. Equality is attained at $(10,0)$, so the maximum is exactly 20. The inequality proof and feasible witness establish both sides.

A reported optimum without a feasible witness or trustworthy bound is weaker. Floating-point tolerances can matter near a covenant. Rational reconstruction or validated interval checks may be appropriate for a small certificate. The economic model should not change its threshold merely to accommodate solver tolerance.

Linear programming does not resolve categorical interpretations, nonlinear prices or integer security units automatically. Those additions can require mixed-integer, nonlinear or hybrid methods. The profile should identify which relations remain outside the exact solver and how their unresolved regions are retained.

\section{Projection can preserve a useful small relation}
Suppose receivables satisfy $A_1=A_0+S-K-W$, with nonnegative write-offs $W$, collections $K$ and known credit sales $S$. If only $A_0,A_1,S$ matter to a future query, eliminating $K$ and $W$ leaves $K+W=A_0+S-A_1$. That sum is determined while its allocation is not.

A credit-loss stress depends on $W$ or surviving customer exposure; a pure cash-history query depends on $K$. Eliminating both into an unnamed residual can make neither query answerable. A better compressed record retains the identified sum and the feasible allocation interval.

If $A_0=30$, $S=40$ and $A_1=20$, then $K+W=50$. With a note limiting write-offs to at most 5, collections lie in $[45,50]$. A cash-flow disclosure of collections 48 then identifies write-offs 2 within this simplified bridge. Each additional source narrows the joint relation.

If currency movements or acquisitions are later admitted, the relation changes. A cached inference $W=2$ must be invalidated unless those terms are established zero. This is an example of model expansion changing an earlier identified quantity. The system should preserve the old inference with its premises rather than silently rewriting its meaning.

\section{Intervals and dependence}
Interval arithmetic can bound expressions quickly, but repeated variables can cause overestimation. If $x\in[0,1]$, the expression $x-x$ is exactly zero. Treating its two occurrences independently gives interval $[-1,1]$. The dependency loss is entirely computational.

For the financial expression $C=95-I$ with $I\in[5,7]$, interval evaluation correctly gives $C\in[88,90]$. But evaluating $C+I$ from separate intervals gives $[93,97]$ instead of exact 95. A relational representation or symbolic substitution preserves the identity.

Overestimation can be acceptable for an outer bound if it is labeled and used with the correct direction. It can become useless near a threshold or create spurious defaults. Refinement can retain shared symbols, affine forms or polyhedral constraints rather than independently widening every intermediate amount.

Nonlinear functions add further issues. A naive interval for $x(1-x)$ on $[0,1]$ is $[0,1]$, while the exact range is $[0,1/4]$. Completing the square, $x(1-x)=1/4-(x-1/2)^2$, identifies the maximum. A small analytic reformulation can be more valuable than a large random sample.

\section{Piecewise contract payoffs}
A floating coupon with floor $l$ and cap $u$ can be written $N\delta\min(u,\max(l,r+s))$ under a stated schedule. Unknown reference rate $r$ partitions into floor, interior and cap regions. The function is continuous but has slope changes.

Suppose principal is 100, year fraction one, spread zero, floor 3 per cent and cap 6 per cent. For $r\in[0.02,0.08]$, interest ranges from 3 to 6. A linear approximation using the midpoint rate five per cent and an unconstrained slope predicts 2 to 8, missing both contractual limits.

If the query is only the interest range, monotonicity gives exact endpoints. If the query includes a refinancing choice triggered by interest above 5.5, a split at $r=0.055$ is needed in the interior region. If the cap changes after a covenant event, contract state must retain that dependence.

A piecewise formula can therefore represent many paths compactly. The implementation needs complete and disjoint guards, correct boundary inclusion and a record of the clause version. A missing cap is an economic-model error; an incorrect inequality at the cap boundary is an implementation error. Their repairs differ.

\section{Nonmonotone outcomes}
Many financial outputs are not monotone in a primitive shock. A rate increase raises floating interest but can increase income on cash, reduce a fixed-rate debt's market value or trigger a favorable option exercise. A customer demand increase can improve revenue while creating a working-capital funding need.

Consider cash after a profitable sale with collection delay. Producing $q$ units costs $4q$ immediately, while revenue $6q$ is collected later. Current cash falls as $q$ rises even though eventual profit rises. A solver maximizing profit may select a quantity that violates interim liquidity.

If financing capacity is $B$ and opening cash $C_0$, feasible production satisfies $4q\leq C_0+B$ before collection. The eventual profit objective $2q$ does not remove this constraint. Aggregating the period so that revenue and cost net immediately changes the feasible set.

A query-driven method must therefore include intermediate constraints and not rely on a generic sign assumption. Monotonicity can be proved locally within a regime and fail when a guard changes. The method should partition or return a qualified bound rather than extrapolate the local result globally.

\section{Integer quantities and minimum lots}
Shares, bond lots and physical units may be discrete. If a company needs net cash 10 and shares sell for 3 with no fees, a continuous calculation requires $10/3$ shares. If only whole shares can be issued, four shares raise 12. The excess cash and dilution are real consequences of the lot constraint.

A minimum issuance size can make a small funding need expensive or infeasible. A fixed fee can make splitting one issue into two change net proceeds. These effects break the proportional invariance that holds for a simple linear loan.

The model should specify whether quantities are continuous approximations or exact units. For large public share counts, a continuous approximation may be immaterial for some queries. For a closely held company or indivisible asset, it may dominate the decision. The approximation needs a bound tied to the question.

\section{Branch-and-bound interpretation}
A mixed discrete-continuous family can be explored through bounds on regions. If a region's best possible result cannot improve the incumbent for an optimization query, it can be pruned under the method's correctness conditions. If a region is infeasible, it can be removed with evidence.

Pruning for computational convenience is different. A low-probability region cannot be removed from a worst-case query merely because it is unlikely. A region far from current data may still be admissible if the observation model permits it. The pruning reason belongs in the result record.

For a safety query, an outer lower cash bound above the requirement can close a region as safe. A feasible witness below it establishes possible failure. Regions whose bounds straddle the threshold remain unresolved and receive refinement priority. This approach can concentrate work where it changes the answer.

The procedure may still face exponential growth. Resource exhaustion should preserve the incumbent witnesses, certified bounds and unresolved regions. Reporting only the best found point hides the remaining uncertainty and can reverse the meaning of a worst-case result.

\section{Sensitivity of identified sets}
A small source revision can change a feasible family discontinuously when constraints meet at a boundary. Suppose exact observations require $x=1$ and an interpretation requires $x\leq1$. The family is nonempty. Revising the observation to $x=1.001$ makes it empty if no tolerance or other interpretation is admitted.

This does not mean the company changed abruptly. It may reveal that exact equality was too strong for a rounded report. Observation precision and numerical tolerance should be audited before interpreting a discontinuity economically.

Conversely, a small genuine change can cross a contractual threshold and have a large effect. Smoothing it away would hide a real mechanism. The analyst must distinguish measurement uncertainty from contractual discontinuity. Both can occur near the same number but require different treatment.

A useful sensitivity report varies source precision, interpretation choices and economic shocks separately. It identifies whether a conclusion is driven by evidence, modeling assumptions or a threshold in the contract. This is more informative than a single generic tolerance parameter.

\section{A bounded implementation profile}
An initial exact profile can support rational linear constraints, finite interpretation choices and piecewise affine event rules. It can return feasible witnesses, exact or certified bounds, and guarded outcome formulas. Nonlinear valuation and unsupported clauses remain explicit extensions.

The profile should reject incompatible units, unbound variables and missing event premises before solving. It should preserve unknown amounts as variables with justified domains. Arbitrary large bounds inserted to make a solver finite should be labeled as computational assumptions and tested for sensitivity.

A later nonlinear profile can use interval methods, constrained optimization or sampling, but must retain the original query target and distinguish local from global results. A local minimum of cash is not a global worst case. A successful solver exit code is not a proof of completeness.

The practical advantage of a bounded profile is that its claims can be checked. It does not redefine the whole mission to exclude difficult contracts. Instead, it provides a reliable core and a precise account of which additional methods are needed for broader company simulation.

\section{Sharing a rate family without sharing away a borrowing right}
\label{sim:continuous-guard-reference}

Return to the capped floating coupon, but now ask whether the company can pay it. The company starts with cash of 5.5, other assets of 94.5 and debt of 100. Its equity is initially zero. One admissible history has a fixed coupon of 5; another has the floating coupon
\begin{equation}
 I(r)=100\min\{0.06,\max\{0.03,r\}\},
 \qquad r\in[0.02,0.08].
\end{equation}
At the earlier observed rate of five per cent, both histories report interest of 5. A later rate of six per cent separates them: fixed interest stays at 5 while floating interest reaches 6. A journal history that retains only the earlier interest amount cannot answer the new payment question.

Each history also has two possible contractual interpretations. A committed facility may be available, or no facility may be available. In this stipulated example the facility lends 1, charges an immediate fee of 0.25 and can be drawn only when interest exceeds 5.5. These terms define a small test case; they are not inferred from the earlier statements. The four combinations of coupon type and facility availability remain separate possibilities. No probability is attached to them.

There is one decision date. Interest accrues before the decision, and the manager submits either a payment order or a combined draw-and-pay order. These are indivisible instructions. An instruction that fails a prerequisite is refused. Contingent instructions, such as \emph{draw only if the bank permits it}, would enlarge the action set and require a fresh analysis.

\subsection{The amount and the right determine different conditions}

A payment order requires enough existing cash. A draw-and-pay order requires the contractual right, the strict trigger and enough cash after the fee. Write $b=1$ when the facility is committed and $b=0$ otherwise. Feasibility is therefore
\begin{align}
 \text{pay:}\quad &5.5-I\geq0,\\
 \text{draw and pay:}\quad &b=1,\quad I>5.5,
                    \quad 5.5+1-0.25-I\geq0.
\end{align}
Because interest never exceeds 6, an eligible draw leaves at least 0.25 cash. At interest exactly 5.5, ordinary payment remains possible and the draw is forbidden. Replacing the strict trigger with a weak inequality would change the contract at that point. At interest 6, the two facility interpretations have identical interest amounts but different feasible actions: the committed company can draw, while the company without the facility can execute neither order.

Accrual is independent of that execution result. The company first debits interest expense and credits interest payable by $I$. If payment refuses, cash remains 5.5 and interest payable remains $I$; equity has fallen to $-I$. A failed payment does not erase the expense or the creditor's claim. If payment succeeds, it removes the payable and cash becomes $5.5-I$. If the draw succeeds, borrowing adds 1 of cash and debt, the fee reduces cash and equity by 0.25, and payment removes the payable. The resulting balance sheet satisfies
\begin{equation}
 \underbrace{(6.25-I)+94.5}_{\text{assets}}
 =\underbrace{101}_{\text{debt}}
  +\underbrace{(-I-0.25)}_{\text{equity}}.
\end{equation}
These journal rules are supplied premises for the example. They demonstrate how refusal and contractual memory enter the comparison; they do not determine a reporting treatment for an actual instrument.

\subsection{A continuous comparison can be exhaustive}

For floating debt, the coupon changes formula at rates 0.03 and 0.06, and the payment trigger changes at 0.055. Suppose also that a coarse signal tells the manager whether the rate is at most 0.04. The six points
\begin{equation}
 0.02,\quad0.03,\quad0.04,\quad0.055,\quad0.06,\quad0.08
\end{equation}
then delimit every relevant change. Keep each point as a singleton and each interval between successive points as an open interval. The eleven resulting regions are disjoint and cover the entire closed rate domain. The fixed-coupon alternative can use fewer regions because it has no rate-dependent coupon or payment trigger.

Within each open region, every financial amount is affine and every signal and eligibility condition is constant. Two implementations therefore agree throughout that region when their affine coefficients and discrete results agree. At singleton regions, compare their values directly. If one implementation uses a coarser partition, compare both on a common refinement and establish that each refined region lies entirely inside one region of each implementation. Sampling a midpoint can identify a known affine piece, but it cannot establish that the list of pieces is complete.

The checked reference makes this comparison for the full journal, cash flows, balance sheet, payable, facility right, observations and action outcomes. Its independent comparator explicitly enumerates all coupon types, facility interpretations and the eleven rate regions. The smaller representation retains 32 mode-specific regions and shares 14 identical calculations. Every original mode and region still points to its calculation. These counts describe this example: sharing reduces duplicated results here, but supplies no general bound on memory use or runtime.

Even this small construction exposed a consequential implementation mistake. Solving the draw-liquidity equation using the interior formula $I=100r$ gives a root at $r=0.0625$. That formula applies only below the cap at 0.06. The root is irrelevant to the capped region and must not be inserted as though the interior formula still governed there. A root finder must check the domain of the formula that generated its candidate boundary.

Open regions also require careful treatment of extrema. On $0.055<r<0.06$, cash after drawing and paying is $6.25-100r$, so
\begin{equation}
 \inf C=0.25,\qquad \sup C=0.75.
\end{equation}
Neither endpoint is attained within that open region. The point $r=0.06$ attains cash 0.25, while $r=0.055$ belongs to the ordinary-payment region. Reporting an interval without its endpoint status loses a distinction needed for strict contractual conditions.

\subsection{Information can prevent an otherwise affordable action}

For the moment retain only histories with a committed facility. Every detailed history admits one of the two orders. This does not imply that the manager can select a feasible order using the information available at the decision date.

Under the coarse signal, \emph{high} means only that $r>0.04$. It includes floating rates with interest below 5.5, where payment is possible but drawing is forbidden, and rates above 0.055, where drawing is possible but ordinary payment fails. The manager must issue the same order at indistinguishable histories. Neither order works throughout this signal class.

Now suppose a second, explicit pre-decision signal reveals whether interest exceeds 5.5. The manager can pay when the trigger is false and draw and pay when it is true. This rule is feasible throughout the committed-facility family. Its feasibility comes from the additional information, not from a better numerical approximation. If unavailable facilities are restored to the family, some high-rate floating histories admit neither order. Even complete information cannot create borrowing capacity.

The exact condition is elementary in this finite action setting. Let $H(o)$ contain all retained histories with observation $o$, and let $A(h)$ be the actions feasible at history $h$. A deterministic policy feasible throughout the family exists exactly when
\begin{equation}
 \bigcap_{h\in H(o)}A(h)\ne\varnothing
 \quad\text{for every observed class }H(o).
\end{equation}
Necessity follows because one observation selects one action, which must work at every compatible history. For sufficiency, choose an action from each nonempty intersection. The finite collection of signals then defines a total policy. This argument depends on the supplied observation map, two atomic orders and robust feasibility across the family. It is not a conclusion about contingent instructions or a probability-weighted objective.

\subsection{What remains safe when calculation is unfinished}

A computation budget can stop evaluation of individual regions without removing them from the family. In the reference, the partition is constructed first and transitions are evaluated when requested. If only one region has been evaluated, the others remain explicitly unresolved. The coupon still has the valid outer bound $3\leq I\leq6$, but that bound alone does not identify a feasible payment policy. Eager construction of the partition itself can become expensive when more continuous variables or contractual alternatives are added.

Approximation needs a separate justification. Suppose an independently established bound gives $|I-\widehat I|\leq\varepsilon$ throughout a region. A strict trigger $I>K$ is certainly true when $\widehat I-\varepsilon>K$, and certainly false when $\widehat I+\varepsilon\leq K$. In the remaining band the bound permits both answers, so the region must be refined or left unresolved. The conclusion follows by comparing the entire allowed interval with $K$; sampled errors do not supply its uniform premise.

Observation thresholds need their own error bounds. Preserving the interest trigger while changing which rate signal the manager sees can still change the feasible-policy set. A reused calculation must therefore retain its rate domain, contract interpretation, clause version, horizon, query, observation map, action class and precision. Changing any of these requires reconsidering the grouping. The example provides an exact one-rate, one-date reference for that discipline. Persistent hidden terms across several dates, higher-dimensional regions, nonlinear guards and defensible approximation bounds remain further method problems.\footnote{The derivation and independently authored comparison are retained in \path{data/accounting-master/continuous-guards-20261026/}. The accepted result is in \path{attempt-02/}; the earlier failed partition and its repair remain available separately.}


\section{Repayment changes tomorrow's feasible payments}
\label{sim:two-date-continuous-memory}

The previous payment calculation ends as soon as interest has been settled. A facility used to make that payment may have to be repaid before the next coupon. To see why the horizon matters, carry the same company forward for one more date. The fixed or floating coupon interpretation persists, as does the existence or absence of the facility. The two reference rates may vary independently within the earlier interval. This rectangular rate family is an assumption for the example, not a probability model.

The facility permits one draw of 1 during the two-date horizon and charges the same immediate fee of 0.25. A first-date draw must be repaid on the second date, after an operating cash receipt and before the second coupon. Repayment extinguishes the facility debt but does not restore the right to draw. The original principal of 100 matures beyond the horizon. We omit tax and interest on the facility to isolate the timing of cash and the persistence of the drawing restriction.

Let $I_1$ and $I_2$ denote the two coupons and $R$ the operating cash receipt on the second date. In the fixed interpretation both coupons are 5. In the floating interpretation each is calculated from its own reference rate using the earlier floor and cap. Receiving one coupon of 5 does not by itself reveal the debt type: floating debt can produce the same amount. Nor does it reveal whether the facility exists.

\subsection{The first order leaves a different second opening}

Ordinary payment on the first date leaves cash $5.5-I_1$. Drawing and paying leaves $6.25-I_1$ and facility debt of 1. On the second date the operating receipt increases cash and income by $R$, after which any outstanding facility debt is repaid. Cash available before paying the second coupon is therefore
\begin{equation}
 C_2=
 \begin{cases}
  5.5+R-I_1,&\text{after ordinary first-date payment},\\
  5.25+R-I_1,&\text{after a first-date draw}.
 \end{cases}
\end{equation}
The second line follows by adding $R$ to $6.25-I_1$ and subtracting the repayment of 1. Although the facility improved first-date liquidity, its fee leaves less cash available over the complete two-date sequence.

Compare two supplied receipt amounts, $R=6$ and $R=7$. When both coupons are 6, the first payment requires a draw and leaves cash 0.25. The consequences at the second date are:
\begin{center}
\begin{tabular}{rrrr}
\toprule
Operating receipt & Cash after repayment & Coupon due & Payment result\\
\midrule
6 & 5.25 & 6 & Refused; payable 6 remains\\
7 & 6.25 & 6 & Paid; cash 0.25 remains\\
\bottomrule
\end{tabular}
\end{center}
The receipt of 6 cannot fund the second coupon. Omitting repayment would incorrectly show 6.25 cash. Allowing another draw would also reverse the answer: the extra 1 less its 0.25 fee exactly fills the 0.75 shortfall. Both mistakes erase a consequence of the first decision.

Before either payment order, the second coupon accrues as interest expense and interest payable. A refused order leaves that payable and the earlier receipt, repayment and fee entries intact. If a first-date order already refuses, this particular reference returns that incomplete prefix without inventing second-date states. Modeling later default, recovery or renegotiation requires additional events.

\subsection{The relevant boundary now joins both rates}

Following a first-date draw, ordinary second payment requires $C_2-I_2\geq0$. Substituting the cash expression gives
\begin{equation}
 I_1+I_2\leq5.25+R.
\end{equation}
For a receipt of 6, the boundary is $I_1+I_2=11.25$. Within the interior floating-coupon pieces it becomes a diagonal boundary in the two rates, $100r_1+100r_2=11.25$. Separate bounds on the two marginal rates do not describe which joint histories cross it. Equality permits payment; larger coupon totals cause refusal.

For a receipt of 7, the corresponding bound is 12.25. Each coupon is at most 6, so their sum is at most 12 and ordinary second payment is feasible throughout the first-draw region. After ordinary first payment, $I_1\leq5.5$; even the receipt of 6 then supplies at least 6 before the next coupon. The failure with the smaller receipt is thus concentrated in histories that needed an earlier draw and then face a sufficiently large combined coupon.

The exact comparison first takes the products of the earlier point and open rate intervals. It then splits those regions at the later linear cash conditions. Every resulting balance and journal amount is affine in $(r_1,r_2)$, while each action condition has a fixed truth value. The independent financial calculation uses the closed-form cash expressions above. It checks that accepted regions cover exactly the feasible domains, that refusals cover the remainder, and that both calculations produce the same complete journal and financial state.

The geometric calculation also retains strict boundaries. On a bounded region described by rational linear inequalities, its closed relaxation is generated by finitely many vertices. The average of all those vertices satisfies a required strict inequality whenever at least one vertex satisfies it strictly: every other vertex has nonpositive slack, so their average remains strictly negative. If all vertices lie on its boundary, no point of the closed region can satisfy that inequality strictly. Applying this reasoning to every strict constraint decides whether the original region contains a point. Extremum attainment is checked by adding the optimizing equality and testing the resulting face in the same way. These arguments apply to the supplied bounded linear regions; a nonlinear or unbounded problem needs another method. The financial comparator shares these geometric primitives, so its agreement does not constitute an independent solver validation.

\subsection{More cash does not reveal a hidden contract}

With a committed facility and a receipt of 7, the refined first-date signal supports the same initial rule as before: pay when the trigger is false and draw when it is true. On the second date the manager observes whether the coupon exceeds 5.5 and whether available cash covers it, together with the earlier signal and chosen order. Ordinary payment is feasible throughout this family.

Some histories also permit a second-date draw: the first coupon was paid without drawing, the facility remains unused and the second coupon exceeds 5.5. There are two such reachable observation classes in this example, and either ordinary payment or drawing is feasible in each. Consequently four policies satisfy the stipulated feasibility conditions. This count does not rank their value. Optional borrowing incurs the fee and leaves debt beyond the current horizon.

Under the coarse first-date signal, the manager still cannot distinguish histories requiring opposite first orders. Raising the later receipt does not repair that information problem. Retaining histories without a facility also leaves some first payments unfunded. A correct continuation therefore preserves contractual memory, cash timing and the information available to the actor together. Renewing the facility, redrawing the debt type or revealing an unobserved trigger would change the problem being solved.\footnote{Complete symbolic states, accepted-domain checks, geometric diagnostics and the retained journal-identity repairs are in \path{data/accounting-master/two-date-continuous-memory-20261026/}. The accepted run is \path{attempt-04/}. This is an exact small reference; demand-driven two-date compression, checkpoint/resume guarantees, source revision invalidation and integration with the common typed executor remain separate work.}


\section{Connecting the history to accounting execution}
\label{sim:guarded-memory-execution}

A financial path can be correct in a small mathematical model and still be translated incorrectly into a company's accounting system. The preceding example treats equity as a single amount after income and expenses. An accounting book usually distinguishes capital, operating income, interest expense and fees. Those distinctions matter when the simulator must produce an income statement, cash-flow statement and retained claim history as well as a closing balance sheet.

The checked connection uses the existing common accounting executor with explicit definitions of the selected events. It does not put the earlier Python posting calculation inside that executor. The comparison instead asks whether independently specified methods produce the same full accounts, journals, cash-flow origins and contractual memory. Two different account charts test whether a method depends on accounting roles rather than convenient account names.

\subsection{Equity agrees only after the accounts are mapped}

Let $B_k$ be an account's signed balance, positive for a debit and negative for a credit. Assets and expenses then normally have positive balances, while liabilities, capital and income normally have negative balances. In this example the reduced equity $E$ used in the earlier calculation corresponds to
\begin{equation}
 E=-B_{\mathrm{capital}}-B_{\mathrm{income}}
   -B_{\mathrm{interest\ expense}}-B_{\mathrm{fee\ expense}}.
\end{equation}
For example, an operating receipt of 7 gives the income account a credit balance of $-7$. Two coupons of 6 give interest expense a debit balance of 12, and a facility fee adds expense of 0.25. With no opening capital, the reported equity after profit is $7-12-0.25=-5.25$. This equals the reduced model's equity without hiding the income statement inside the capital account.

An interest accrual $I$ adds $I$ to expense and $-I$ to interest payable. Under the mapping, equity falls by $I$ and the creditor's claim rises by $I$. Paying the coupon reduces cash and removes the payable, leaving profit unchanged. A combined draw-and-pay instruction has the following signed account changes:
\begin{align}
 \Delta B_{\mathrm{cash}}&=0.75-I,&
 \Delta B_{\mathrm{facility\ debt}}&=-1,\\
 \Delta B_{\mathrm{fee\ expense}}&=0.25,&
 \Delta B_{\mathrm{interest\ payable}}&=I.
\end{align}
Their sum is zero. The cash-flow record separately retains borrowing of 1, the fee payment of 0.25 and the interest payment of $I$. Netting those amounts in the journal's cash entry must not erase their origins in the cash-flow statement.

Repayment removes the facility liability and reduces cash by the same amount. Its accounting balance becomes zero, but the separate memory that the facility has already been used remains true. Each method also checks that interest payable agrees with accrued coupons less paid coupons, that the event occurs at the declared phase and date, and that its prerequisites are satisfied. A refused payment leaves the accepted prefix unchanged, including the coupon already accrued.

For the concrete execution comparison, the two decisions are assigned to 30 June and 31 December 2027 within one supplied reporting period. The coupon amounts remain the one-period charges of the earlier example. These dates test ordering; they do not supply an actual instrument's day-count convention. Each comparison checks the complete account balances, income statement, balance sheet, gross cash flows, journals and memory. Applying the equity mapping also recovers the earlier independent financial calculation.

\subsection{A checked path is not a checked continuous region}

The common executor accepts concrete rational amounts. The selected comparison therefore uses exact points at the coupon cap, at the first trigger, and at and immediately to either side of the later joint cash boundary. Those executions exercise important failure mechanisms. They do not by themselves establish that every rate in a continuous region produces the correct accounting result.

Within a region where the guards are fixed, the amount relations above are affine. If the opening accounts map correctly, each event has the stated affine effect, and the same guard, timing and memory conditions govern both descriptions, induction over the accepted events gives the same mapped financial quantities. The conditions carry the argument: a wrong guard or lost drawing right can change which event occurs even when its monetary formula is correct. A general checker must verify those conditions and every supported method, rather than replace that obligation with a larger sample of rates.

The concrete evidence at this stage has two parts. The exact continuous reference establishes its financial formulas and feasible domains. The concrete common-engine comparison checks complete accounting execution at declared points. The explicit mapping and local event derivation motivate the selected symbolic certificates in Section~\ref{sim:certified-rate-regions}; extending them to general methods remains a separate problem. Source-backed terms and actual accounting applicability also remain distinct from agreement between two synthetic representations.

\subsection{Stopping computation must preserve what remains unknown}

Suppose the simulator has executed the first coupon and then reaches its computation budget. A useful saved result contains the complete accepted book and journal, the next command, and the commands still pending. It must also retain the alternative histories and regions that have not been evaluated. A successful payment on one path supplies no reason to delete a different path that may fail tomorrow.

The selected checkpoint reference validates a saved concrete prefix by replaying it from its declared opening and method registry. It compares the complete transition records and current state before continuing. The tested interruptions recover the same result as uninterrupted execution, including the first-date and second-date refusal prefixes. Altering cash while keeping the book balanced, changing paid-coupon memory or removing the used-facility flag fails that comparison. Merely recomputing a checksum on the altered record would not establish that the altered history followed from its opening.

A zero budget attempts no new financial event. It can still incur the cost of checking the saved prefix. More generally, the budget in this reference counts new event attempts, while replay validation costs additional work. A practical implementation would need a separately justified way to reduce that cost. The present result establishes a small exact recovery comparison, not an efficient or crash-durable storage system.

The queue also distinguishes a completed concrete program from a resolved continuous region. For the two-date example, the declared roster contains four persistent contract combinations, $11\times11$ rate-atom pairs, four order pairs and two receipt hypotheses:
\begin{equation}
 4\times121\times4\times2=3{,}872
 \quad\text{region and instruction obligations}.
\end{equation}
A finite witness cannot mark any of those whole regions complete. Their identities remain unresolved under this concrete checkpoint even when individual concrete jobs succeed or refuse; Section~\ref{sim:certified-rate-regions} supplies separate certificates for the selected whole regions. Within the finite job queue, pending, completed and refused jobs form an exact partition. A refused job retains both its payable and its unexecuted commands; it is a completed assessment of that particular instruction sequence, not proof that every alternative has failed.

Finally, reuse depends on what was asked and what information was available. The saved record names the source revision, clause, actor's observation map, available actions, horizon, query, precision and method. A changed scope refuses reuse under the old record. An earlier query can still use its earlier explicitly supplied scope, preserving the distinction between later knowledge and information available at the time. Those bindings do not authenticate the source or establish when a real manager received it. Source admission, extension beyond the selected symbolic domains, efficient refinement and durable recovery require their own evidence.\footnote{The complete common-engine and checkpoint comparisons are retained in \path{data/accounting-master/guarded-memory-runtime-20261026/}, respectively under \path{attempt-01/} and \path{checkpoint-attempt-01/}. The former checks concrete accounting execution; the latter checks replay and queue preservation. Neither promotes a concrete witness to a whole-region certificate.}


\section{Proving a result for an entire rate region}
\label{sim:certified-rate-regions}

The second coupon gives a useful test of what a simulation has actually established. Suppose the company drew its once-only facility at the first date. After the operating receipt and repayment, cash available for the second coupon is $5.25+R-I_1$. A payment succeeds exactly when
\begin{equation}
 I_1+I_2\leq 5.25+R.
 \label{eq:certified-second-payment}
\end{equation}
At $I_1=I_2=6$, a receipt of 6 leaves only 5.25 before the payment, so the coupon remains unpaid. A receipt of 7 leaves 6.25 and permits payment, with 0.25 left over. A simulator that happens to test only the latter case has learned nothing about the former. Even a large set of successful sample executions does not exclude an omitted boundary or an incorrect condition between the samples.

The selected example now has a stronger, exact connection to the declared accounting methods. It proceeds in three stages: establish each method's equations, substitute those equations through the dated sequence, and divide the rate domain wherever a condition changes. This division concerns computational work. It does not assert that a company, a manager or an observer learns which historical explanation is true.

\subsection{Compare the method with an equation, not another sample}

An affine monetary expression is a constant plus a weighted sum of quantities. For this example, those quantities include account balances, accrued and paid coupons, event amounts and the supplied parameters $I_1,I_2,R$. Write an expression as
\begin{equation}
 f(x)=c+\sum_j a_jx_j,
 \label{eq:certified-affine-form}
\end{equation}
where each coefficient and the constant are exact rationals. The accounting language expresses these calculations through constants, references, addition and subtraction. A constant contributes to $c$; a reference contributes one coefficient; addition and subtraction combine the coefficient lists. Induction on these four constructions shows that the resulting list represents precisely the original expression. Two identical lists therefore give the same amount for every known typed valuation, rather than merely for the values tested in a run.

For instance, the combined borrowing and payment changes cash by $3/4-I_t$. The corresponding coefficient of the due coupon is $-1$ and the constant is $3/4$. Doubling both interest expense and the credit to interest payable would still balance the journal. It would change these independently specified account relations and fail the comparison. Balance is a necessary accounting constraint, but it cannot establish that the right transaction was recorded.

The comparison also preserves the conditions on the transaction. A numerical comparison is represented by its exact left-minus-right affine expression and its comparison sign. The facility trigger $I_t>11/2$ retains a strict sign; replacing it by $I_t\geq11/2$ fails even if a chosen set of sample rates misses equality. Boolean conditions retain their operands, and memory assignments retain their source and destination. The same check covers cash-flow origins: a borrowing receipt, fee and interest payment cannot be replaced by one net cash movement.

Parameters enter through a separate selected template. Its coupon and receipt amounts, contract mode and facility right become ordinary typed literals before execution by the common runtime. Exact placement checks bind the parameters to the opening state, event inputs and selected-amount conditions. The surrounding method records retain their account maps, units, dates, schemas, declared dependencies and refusal policies. This is a certificate for the inspected selected profile. It does not extend the runtime to arbitrary symbolic programs or prove an unrestricted compiler correct.

\subsection{A correct event still needs a correct predecessor}

Local agreement leaves a temporal question. When a method reads the payable or the facility balance, does it receive the state produced by the preceding accepted event? Substitution answers that question for the selected sequence. Every posting and memory assignment reads the same pre-event state; the resulting changes become visible together. The next event receives the complete new state, including accrued and paid coupon history.

After a first-date draw and payment, cash is $25/4-I_1$. The receipt adds $R$, and repayment subtracts the borrowed unit. Substituting successively gives
\begin{align}
 C_{1+}&=\frac{25}{4}-I_1,\\
 C_{\mathrm{receipt}}&=\frac{25}{4}-I_1+R,\\
 C_{2-}&=\frac{21}{4}-I_1+R.
\end{align}
The second accrual changes expense and the payable, leaving this cash unchanged. Substituting $C_{2-}$ into the payment condition $I_2\leq C_{2-}$ yields equation~\eqref{eq:certified-second-payment}. Repayment removes the facility liability while leaving the used-facility flag true. Omitting that memory would admit a second draw which the contract in this example forbids.

The checked substitution compares all signed accounts, memory fields, gross flows and conditions with independent complete formulas. It includes both refusal positions. A first refusal leaves the first coupon accrued; a second refusal retains the accepted first-date events, receipt, repayment and second accrual. It does not replace the final payable by zero merely because the proposed payment appeared in the command sequence.

At this stage a branch can be written down even when no rate makes it feasible. For example, the formal second-draw branch can be constructed after a first draw, but its unused-facility condition is false. Distinguishing an algebraically described branch from a possible history is the task of the next stage.

\subsection{Why an empty branch needs a proof}

Within each rate atom the capped coupon is constant or linear in the corresponding rate. A point atom uses its exact constant coupon. Substituting those pieces into the conditions gives inequalities in two variables, $r_1$ and $r_2$. Each comparison divides the current region into the points satisfying it and the points satisfying its exact negation. For $f\leq0$, the other child is $-f<0$; for $f<0$, it is $-f\leq0$. Equality belongs to one child, never to both and never to neither.

A discarded child carries a short arithmetic reason why its inequalities cannot hold together. Suppose the child requires
\begin{equation}
 a_i^{\mathsf T}r+c_i\leq0,\qquad i=1,\ldots,m,
 \label{eq:certified-region-constraints}
\end{equation}
with some signs possibly strict. Choose rational weights $\lambda_i\geq0$ that cancel both variable coefficients:
\begin{equation}
 \sum_i\lambda_i a_i=0.
\end{equation}
Adding the weighted inequalities then requires $\sum_i\lambda_i c_i\leq0$. If that constant is positive, the requirement is impossible. A zero constant also proves impossibility when at least one strict inequality receives positive weight: the sum would require $0<0$. These implications follow directly by multiplying and adding the original inequalities. The verifier checks the weights against those original rows rather than trusting a solver's declaration that a branch is empty.

The distinction matters even in one dimension. The pair $x<0$ and $-x\leq0$ has no solution: equal unit weights produce $0<0$. Weakening the first inequality to $x\leq0$ admits $x=0$. A certificate that discards the strictness would incorrectly remove that boundary point.

The reference producer eliminates the two variables with exact rational arithmetic and retains the weighted combinations it uses. A separate checker verifies each proposed contradiction. A retained leaf instead carries a point satisfying every inequality, including every strict sign. That point proves the leaf is nonempty. The affine method and sequence argument, together with the predicates along the branch, establishes the financial relation throughout the leaf; the witness alone does not.

\subsection{What can now leave the unfinished-work list}

For the declared two-date model, all $3{,}872$ rate-region and instruction obligations have a checked partition. The partition contains 918 nonempty accepted leaves and 2,960 nonempty refusal leaves; 4,772 additional children are discarded with checked contradiction certificates. These are computational counts, not probabilities. They depend on the selected atom refinement and cannot be interpreted as the likelihood of repayment or default.

The saved proof archive is re-read before its jobs enter a new certified queue. Each entry binds the source revision, clause, method and sequence certificates, domain, query and horizon. A pending job leaves the queue only with the corresponding verified proof. Zero-budget and interrupted JSON continuations retain the same final partition as uninterrupted processing. Checking the saved proof registry costs work in addition to the budget for newly completed jobs.

The earlier concrete checkpoint correctly kept these regions unresolved under its own evidence. The new result supplies separate whole-region certificates; it does not reinterpret the earlier witness as a proof. The distinction is useful beyond this example: completing a computation should mean satisfying the evidence obligation attached to that job, not merely reaching its last instruction.

This result concerns known exact values, the supplied coupon and once-only facility terms, two bounded rates and the selected event sequence. It leaves a further decision problem. An actor can choose differently only when the actor's information distinguishes the relevant histories. Consequently, a list of financially feasible region-specific actions is not yet an admissible policy. Section~\ref{sim:certified-observation-policies} binds these leaves to the supplied observation map and checks decisions that must agree across indistinguishable histories. Section~\ref{sim:source-timed-policies} supplies a conditional source-revision and actor-access binding for the selected public-document profile; actual company evidence and broader information structures remain separate requirements. Higher-dimensional and nonlinear relations, unknown-value semantics, actual accounting authority, numerical quantisation and crash-durable recovery require additional evidence.\footnote{The local derivations and immutable comparisons are retained under \path{data/accounting-master/symbolic-method-certificate-20261026/}, \path{symbolic-path-composition-20261026/} and \path{exact-region-certificates-20261026/}. The last directory includes failed diagnostics, the exact partition archive, an independent weighted-sum checker and separate saved-archive/queue validation. No new literature proposition is asserted by these local calculations.}


\section{A feasible action for each world is not yet a policy}
\label{sim:certified-observation-policies}

Return to the first coupon, when cash is 5.5 and the once-only facility is available. Two companies can produce the same coarse observation that the market rate is high. In one, a fixed coupon is 5; in the other, the floating coupon at the cap is 6. Paying is feasible for the first company, while drawing is barred by the strict trigger. The second company cannot pay from existing cash and needs the facility. If the manager sees only the coarse high signal, the two individually feasible instructions do not form a decision rule the manager can follow. The rule asks for information that the manager does not have.

This is the difference between a financially possible action and an information-feasible policy. Let $w$ denote a possible history and $O_t(w)$ the manager's observation history when the date-$t$ instruction is chosen. An admissible policy must satisfy
\begin{equation}
 O_t(w)=O_t(w')\quad\Longrightarrow\quad
 \pi_t\bigl(O_t(w)\bigr)=\pi_t\bigl(O_t(w')\bigr).
 \label{eq:certified-information-policy}
\end{equation}
Here $O_t$ includes remembered earlier observations and the manager's own earlier instructions. The condition gives one decision to each information cell. It does not require that the underlying histories become identical or that the manager eventually learns the contract type.

\subsection{Refining an observation without inventing knowledge}

The selected first coarse signal reports whether $r_1\leq0.04$. The refined signal adds whether the first coupon exceeds 5.5. In the pair above, refinement distinguishes the two required first instructions. It still does not reveal the facility right or generally identify a persistent coupon type. At the second date the manager remembers the first signal and instruction, then observes whether the second coupon crosses the trigger and whether the cash then available covers it. The receipt hypothesis is held fixed for each comparison.

A financial region need not have a single observation label. For example, the condition that drawing plus cash covers a coupon is weaker than the condition that existing cash alone covers it. A region on which a draw is financially possible can therefore straddle the manager's ordinary cash-coverage observation. The checked construction splits that region on the cash-coverage inequality and its exact complement. Empty children carry the contradiction proofs of the preceding section; nonempty children retain their symbolic financial outcomes. Whole-atom proofs likewise establish that the first low/high signal and each coupon-trigger signal are constant where they are used. The observation is thus valid throughout the cell, rather than assigned from one representative rate.

This timing is consequential. The first decision does not receive the second rate or a later cash balance. The second decision uses cash before its proposed payment. Using terminal cash would let the result of an action determine what the manager supposedly observed before choosing it. Including the hidden facility right in the signal would change the information problem, even if every resulting journal still balanced.

\subsection{Which instructions survive every indistinguishable history?}

For each possible first decision rule, retain every history consistent with that rule. A first-date refusal rejects the rule. At each reachable second-date information cell, intersect the allowed instructions across the retained histories. An empty intersection means there is no action that works in every history the manager still considers possible. Otherwise choose one instruction from each cell's intersection. Cells that cannot be reached under the first rule need no arbitrary decision in this comparison.

An independent formulation expresses the same condition as Boolean restrictions. Let $X_o$ say whether the first instruction at signal $o$ is draw, and $Y_h$ the second instruction at remembered history $h$. A certified first-date refusal under instruction $a$ forbids $X_o=a$. A certified second-date refusal under first instruction $a$ and second instruction $b$ forbids their joint selection:
\begin{equation}
 (X_o\ne a)\ \lor\ (Y_h\ne b).
\end{equation}
Enumerating the satisfying assignments and retaining only reachable second-date choices gives the complete policy set. This independently authored constraint calculation agrees with the action-intersection calculation and with the earlier finite-policy reference. Agreement concerns whole policies, not only the number found.

Across the eight supplied problems---receipts 6 or 7, coarse or refined observations, and committed or mixed facility rights---only the receipt-7, refined-observation, committed-right problem has robustly feasible policies. It has four. The first instruction is fixed by the refined trigger. After a first draw, the exhausted facility cannot be used again and the second instruction must be pay. After a first pay, two reachable second-date cells allow either pay or draw, giving $2\times2=4$ reachable policies. The model supplies no preference or objective that ranks them.

The failures have different causes. Coarse information combines incompatible first instructions. A family containing an unavailable facility includes a high-coupon world with no feasible first instruction. Receipt 6 leaves the double-high path short after repayment, despite knowing the trigger and having the original right. Removing such a world or ignoring its refused branch changes the robust question; it does not solve it.

These results concern the explicit family of possible histories and the supplied observation maps. They assign no probabilities to that family, and a feasible robust policy is not a forecast or an optimal policy. Applying the calculation to an actual manager next requires evidence about source revisions, access times and observations. That evidence must preserve what was available when a decision was made rather than replacing it with a later corrected account.\footnote{The complete observation refinements, whole-atom proofs, two policy formulations, source locks and eight complete comparisons are retained in \path{data/accounting-master/certified-observation-policy-20261026/attempt-01/}. The reference checks all 3,872 saved financial jobs, with 1,766 complementary observation refinements. These are exact calculations for the selected synthetic information structure.}


\section{Which source could this actor use at this date?}
\label{sim:source-timed-policies}

The four feasible policies in the preceding example depend on two supplied interpretations: which receipt and facility-right family is possible, and which signals the decision maker observes. A revised document can change either interpretation. It can also leave both interpretations unchanged while correcting the evidence supporting them. The simulator needs to distinguish those cases before it decides which calculation to repeat.

Consider a correction released on 3 January to a document used for a decision on 2 January. The correction may improve a retrospective account of the company, but it cannot be inserted into what the actor knew on 2 January merely because the simulator downloads it later. Conversely, a document may have been available to the actor on 1 January but entered the simulator on 3 January and passed internal review on 4 January. On 4 January the system can reconstruct the earlier information set if the retained evidence supports that access. It has not established that the reconstruction itself was run on 2 January.

These are distinct dates with distinct subjects. Publication concerns the document's release. Access concerns a particular actor. Ingestion and review concern the information system. Effective time concerns the economic period or event to which the interpretation applies. A publication timestamp alone does not establish personal access; a download timestamp does not establish either the original publication or the document's economic effect.

\subsection{Select revisions before selecting the policy problem}

The selected construction uses an immutable record for each supplied interpretation. The record identifies the source bytes, the interpretation's subject, the previous revision it replaces, its effective time, and separate observations of publication, access, ingestion and review. The caller supplies the evidence accepted for these timing observations. Hashes preserve the identity of those bytes and observations; they do not establish their truth.

For a request, retain revisions that are effective for the chosen economic event, available to the actor by the decision time, and ingested and reviewed by the system's reconstruction cutoff. Known later revisions are excluded from that earlier view. A missing relevant time leaves the selection unresolved. Within a subject, select a revision only if it is the unique remaining head of its replacement chain. Two incomparable available corrections preserve a conflict. Retraction of a required interpretation leaves a missing input rather than silently restoring its earlier assertion.

This procedure is deliberately restricted to a declared source scope. An unrelated management note cannot change the coupon interpretation merely because it appears in the same registry. A relevant correction must be inspected again on each request, however. Comparing a cached view identifier without checking for newly available revisions could miss exactly the change that should invalidate the result.

In the reference example, one required subject supplies receipt 6 or 7 and the mixed or committed facility-right family; the other supplies the coarse or refined observation scheme. These choices map to the eight previously certified problems. They preserve the same possible persistent contract types, complete financial rows and rejected actions within the selected family. No convenient prior or likelihood is introduced to make a source selector callable. The question remains whether one policy works in every history the actor cannot distinguish.

The timing checks reuse the existing source and policy-information mechanisms. The selected adapter accepts explicit whole-second UTC instants and public-document evidence with a separate supplied actor-access observation. It refuses finer or incomplete clock representations instead of rounding them into eligibility. Section~\ref{sim:exact-temporal-views} extends this selected reference to exact clock ranges and an explicit private-delivery route; broader information practices remain separate. An unresolved access record is therefore a substantive result about this request, not an instruction to substitute the publication date.

\subsection{Reuse a calculation without reusing an obsolete source binding}

Suppose a replacement document supports exactly the same receipt, facility-right family and observation scheme. The set of robust policies can remain unchanged. Nevertheless, a report that still cites the replaced source would give the wrong explanation of which evidence supported the current request. Numerical equality and evidential identity answer different questions.

Let $F$ denote the retained financial relations and certified rows, $s$ the selected receipt and facility family, and $o$ the observation scheme. Write the complete reachable policy calculation as
\begin{equation}
 \mathcal P(F,s,o).
 \label{eq:source-policy-dependency}
\end{equation}
Its numerical inputs do not include a document's filename or the system's ingestion date. Those records determine whether the interpretation $(s,o)$ is available for the actor and request. The reported result therefore needs both $\mathcal P(F,s,o)$ and the source, revision, actor and cutoff records that justify using these inputs.

The checked dependency relation makes this separation explicit. Certified financial evidence supports the chosen financial scenario; that scenario and the observation scheme support the policy set. Separate source bindings attach the selected financial and observation interpretations to their revision histories. The final actor binding joins the numerical policy set to both source bindings and to the decision, reconstruction and effective cutoffs.

A same-meaning correction changes a source binding and the final actor binding while preserving the numerical policy set. A change from coarse to refined information preserves the financial relations but requires a new observation and policy calculation. A changed receipt or facility family changes the financial scenario and its dependent policy problem. Changing only the actor or cutoff requires a new actor binding even when the selected interpretations and numerical answer happen to be identical. A correction known to arrive after the old cutoff changes none of that earlier view's selected dependencies.

Reuse is justified by those dependencies, not by observing that the new and old policy counts coincide. The reference compares the entire reused policy set and its complete world-row population with a fresh calculation. It also requires the saved inputs and method bytes to match expected identities supplied outside the transported bundle. Removed or duplicated rows, changed hidden rights or signals, altered policy references and rewritten prior results are rejected. Recomputing an inner checksum after changing a saved result does not make it the previously accepted result.

The scope remains conditional. The caller must establish which source and timing evidence is accepted, and a complete history family for an actual company may require many more subjects than these two. The result shows how revision-sensitive selection can coexist with numerical reuse while preserving the robust question. It establishes neither authenticated storage nor crash recovery, and the small exact comparison supplies no performance ranking.\footnote{The immutable source-view cases are retained in \path{data/accounting-master/certified-source-views-20261026/}; the separate dependency and complete cold-versus-reused comparisons are in \path{data/accounting-master/source-policy-invalidation-20261027/}. Both preserve their supplied source scope and withhold actual-source, accounting, real-actor and forecast admission.}


\section{A date range can leave several information views possible}
\label{sim:exact-temporal-views}

A filing marked as published on 2 January does not establish whether a manager had read it at the beginning of that day. Nor does a correction received within a narrow interval around a decision establish whether the original or corrected interpretation was available for that decision. These are questions about the order of events. Assigning midnight or the interval's midpoint would settle them by assumption.

The selected successor retains that uncertainty explicitly. An exact timestamp is a point; a calendar day in UTC is a half-open interval from its midnight to the next midnight. A supplied time range has the same half-open convention. An unknown access time remains unknown. Exact rational seconds preserve differences smaller than a microsecond and equivalent fixed-offset spellings without passing through a floating-point or microsecond clock.

If the access time lies in $[a,b)$, a decision at $b$ is certainly after access. A decision at $a$ can coincide with access or precede it: the lower endpoint belongs to the interval. The distinction disappears if an implementation replaces the interval with a single representative time. It also disappears if a strict upper bound is inadvertently changed to a closed one.

\subsection{Availability depends on the route and on joint chronology}

A public-document retrieval and a private delivery impose different conditions. In the first route the source must have been released publicly before the claimed access. In the second, an actor-addressed delivery can establish access before any public release. The private route neither invents a public timestamp nor makes the document available to other actors. Both routes still require evidence about that actor's access; publication alone cannot supply it.

The earlier whole-second reference did not distinguish these routes. One of its retained cases places actor access before a later public release. Under the new, explicit public-retrieval premise, that chronology is inconsistent. It can be meaningful under a private-delivery premise. Preserving the difference is necessary: changing the access mechanism changes the question rather than improving timestamp precision alone.

Clocks are also dependent. Authorship precedes ingestion, ingestion precedes review, and authorship precedes access. Public retrieval additionally follows publication. The selected revision profile requires nondecreasing authorship, ingestion, review and same-actor access along a replacement chain, while permitting retrospective effective dates. These are declared premises of this bounded source model. A document practice that violates one of them requires another profile; a corrected economic effective date is not forced forward merely to preserve a convenient time ordering.

An unknown clock need not make every conclusion unknown. If a private access time is supplied before the decision and the document was necessarily authored before that access, the chronology can establish the relevant authorship inequality without identifying its timestamp. The result must retain the unknown-time marker. Logical implication from other premises does not turn an unobserved time into an observed fact.

\subsection{Partition clock possibilities without enumerating timestamps}

Let $x_e$ be the time of event $e$, measured in exact UTC seconds on the retained calendar axis. Point and interval observations, chronology and cutoff conditions become difference inequalities. A non-strict edge represents
\begin{equation}
 x_v-x_u\leq c,
 \label{eq:temporal-difference-edge}
\end{equation}
with a strict variant when the right endpoint is excluded. A fixed origin converts an absolute bound into the same form. The source identities, actors, access routes and original precision remain attached to these mathematical restrictions; the inequalities alone do not establish the source facts.

For example, publication before access contributes $x_{\mathrm{publication}}\leq x_{\mathrm{access}}$, while access by the decision contributes $x_{\mathrm{access}}\leq t_{\mathrm{decision}}$. Keeping both inequalities in one system prevents the calculation from accepting a combination that separate per-clock classifications would mistakenly permit.

For each revision, eligibility is a conjunction of its required actor, system and effective-time guards, $g_1,\ldots,g_m$. The complementary cases can be written as a disjoint first-failure partition:
\begin{equation}
 \neg(g_1\land\cdots\land g_m)
 =\neg g_1\;\lor\;(g_1\land\neg g_2)\;\lor\;\cdots\;\lor\;
 (g_1\land\cdots\land g_{m-1}\land\neg g_m).
 \label{eq:temporal-first-failure}
\end{equation}
Each branch preserves the earlier satisfied guards. Repeating the construction for the relevant revisions covers every clock assignment admitted by the original ranges and chronology. Eligibility and the selected revision heads are constant within each surviving branch. The procedure therefore represents families of timestamps symbolically; it does not need an explicit tree point for every possible instant.

An empty branch requires a checked contradiction. Summing original difference inequalities around a cycle cancels all clock variables. A negative right-hand sum is impossible; a zero sum is also impossible if at least one inequality in the sum is strict. A nonempty branch instead carries exact rational clock values satisfying every original inequality. The checker verifies those values directly and reconstructs every complementary branch before accepting the partition.

The witness establishes nonemptiness, while the guards and complete partition establish the scope of the conclusion. A few feasible timestamp examples would establish neither exhaustive coverage nor the absence of an omitted information view. Conversely, reaching the computation budget leaves an unfinished temporal family. The selected reference then reports incompleteness and does not present the already examined views as the full set.

\subsection{Different timing views can support different policies}

Return to receipt 7 and the committed facility. Suppose a correction changes the available observation scheme from coarse to refined, but the actor's access interval straddles the first decision. There are two remaining timing views. In the view where the correction has not yet arrived, the coarse signal leaves incompatible first instructions in the same information cell and no robust policy is feasible. In the view where it has arrived, the refined signal supports the four previously derived policies. The complete financial row population is retained in both calculations.

The timing evidence supplies no probability for either view. Selecting the four-policy answer would discard the other possibility. Reporting the two answers separately is therefore a substantive conclusion. Constructing one decision rule that is admissible when even the observation mechanism is uncertain requires a common specification of what the actor sees and how actions are represented across those mechanisms; a naive intersection of differently labelled policy tables is insufficient.

The exact comparison verifies interval boundaries, equivalent offsets, unknown and untrusted clocks, public and private access, correlated correction times, conflicting revisions and retractions. The saved partitions retain unresolved source outcomes as well as selected ones. Their scope is the supplied chronology and access model, not an empirical statement about how a real company releases information. General revision practices, authenticated evidence and durable recovery remain separate questions; Section~\ref{sim:temporal-family-policies} develops complete-family reuse and common decisions across possible observation mechanisms.\footnote{The derivation, original-edge checker, complete saved temporal partitions, independent finite order comparisons and full policy sets are retained in \path{data/accounting-master/source-time-intervals-20261027/}. The reference preserves its failed namespace diagnostic, the explicit public-route difference from the whole-second profile, and the six-revision and computation-budget limits. No clock probabilities or actual actor-access admission are introduced.}


\section{Preserving possible views does not yet choose a common policy}
\label{sim:temporal-family-policies}

The correction interval leaves two questions. A calculation must first preserve every information view consistent with the supplied clocks. The decision maker must then choose instructions that can actually be followed when the relevant view may itself be unknown. Keeping both policy tables answers the first question. It does not, by itself, answer the second.

\subsection{Share calculations while retaining their different supports}

A complete temporal family carries the original clock restrictions, its checked partition, every nonempty branch and the source outcome on that branch. Several branches can select the same revision; several different revisions can also imply the same financial and observation problem. The numerical calculation can then be stored once, provided that all the distinct source and timing supports remain attached to it. Removing a repeated numerical answer must not remove a possible history.

The selected binding reference therefore compares complete policy sets and their full financial-row populations. A same-meaning correction can reuse the numerical result while changing the source binding. A changed financial assumption or observation scheme requires its corresponding numerical problem. Even two empty policy sets retain their different scenario meanings. Their common cardinality, zero, is not sufficient evidence that the underlying questions are interchangeable.

The request and saved family are checked against expected identities supplied outside the transported record. The request includes the actor, original sources, clock observations, caller's admitted evidence and computation budget. The family retains the full temporal proof and each unresolved outcome. These checks detect alteration and missing support relative to that supplied record. They do not authenticate a document or establish that the actor actually received it.

A budget stop needs special care. The earlier temporal procedure does not save a partial partition that can be resumed. The binding reference therefore preserves the entire original clock domain as unresolved and supplies no completed-view policy answer. This conservative remainder keeps all possibilities, at the cost of repeating work later. Describing it as an incremental frontier would claim a computational capability that has not been implemented.

\subsection{Decisions must agree on the histories the actor can distinguish}

Let $z$ denote a possible observation mechanism and $w$ a possible financial history. At decision date $t$, the actor receives a message $m_t(z,w)$ from that mechanism. The message may differ across histories, and the meaning of a given message may differ across mechanisms. Let $\eta_t(z)$ record a mechanism identifier only when the model explicitly makes it observable by that date. Without such an observation, $\eta_t$ has the same null value for every mechanism.

For the two-date example, the information used by a deterministic policy is
\begin{align}
 h_1(z,w)&=\bigl(m_1(z,w),\eta_1(z)\bigr),\\
 h_2(z,w,a_1)&=\bigl(h_1(z,w),a_1,m_2(z,w,a_1),\eta_2(z)\bigr).
 \label{eq:common-mechanism-history}
\end{align}
The second history retains the first message and the actor's own instruction. It contains no future message, hidden financial state or internal proof identifier. The action at each date is a function of the corresponding history, $a_t=\pi_t(h_t)$. Consequently,
\begin{equation}
 h_t(z,w)=h_t(z',w')\quad\Longrightarrow\quad
 \pi_t\bigl(h_t(z,w)\bigr)=\pi_t\bigl(h_t(z',w')\bigr).
 \label{eq:cross-mechanism-nonanticipation}
\end{equation}
This is the common-policy restriction: indistinguishable available histories require the same instruction. It applies across possible mechanisms as well as across financial histories within one mechanism. Feasibility must then hold for every retained world under the selected stress question. No probability distribution is needed for that requirement.

Internal source hashes can identify $z$ for the analyst without identifying it for the actor. Putting such a hash into $h_1$ would add information and could change the answer. Conversely, if a real message identifies its source or interpretation, omitting that information would solve a different problem. The observation declaration must say which case is intended and why.

\subsection{Four policies in each view can leave no common instruction}

Consider again receipt 7, the committed facility and the refined financial problem. Its complete row population supports four policies. For a concrete test, call the first messages $L$, $H$ and $X$. Mechanism A sends these messages for, respectively, low conditions without a trigger, high conditions without a trigger, and high conditions with a trigger. Mechanism B exchanges the meanings of $H$ and $X$. These are deliberately supplied synthetic encodings; they are not an empirical claim about a company's communications.

Each known encoding preserves the original four policies by relabeling their observations. The difficulty appears when the actor sees $H$ but does not know which encoding applies. Under A, the instruction must be pay; under B, it must be draw. A common policy cannot satisfy both requirements. The two individually feasible problems therefore have no jointly feasible policy when the mechanism is hidden.

If the encoding is identified before the first decision, the two message histories are distinguishable. One of four policies can be selected for each encoding, giving $4\times4=16$ feasible joint policies. If the encoding is identified only at the second date, the earlier contradiction remains. Later learning cannot repair a first instruction that had to satisfy incompatible requirements at the time it was chosen.

The conclusion depends on when the incompatibility arises. In a separate elementary liquidity example, the only first instruction is wait. At the second date, a company with cash 1, an amount due of 1 and no facility must pay. Another possible company has cash 0, the same amount due and a facility of 1, so it must draw. Two mechanisms exchange the meanings of the messages describing those circumstances. Hidden message meaning leaves no common second instruction; revealing the mechanism at the second date restores one feasible policy because waiting had created no earlier conflict. If the second message itself has an unambiguous meaning, learning can resolve the uncertainty without revealing a mechanism's name. This example concerns stipulated cash feasibility, not the accounting recognition of a facility or a complete company report.

A different case shows why late revelation can expand a policy set even when the mechanisms have the same financial meaning. Two possible correction heads can both select the same refined problem. All four of its policies share the same first instructions; only their later choices differ. Learning the source identity at the second date then permits four continuations for each source, giving 16 joint policies. When that identity is never observed, the continuations must agree and only four policies remain. A first attempt incorrectly expected four in both cases; the complete-set comparison exposed the omitted late distinction, and the corrected independent relabeling preserves the failed attempt as evidence.

\pagebreak
\subsection{What a complete conditional answer must retain}

The original coarse/refined correction example remains different from the swapped-encoding example. The coarse member already has no feasible policy. Identifying that member does not supply the refined signal it lacks. The joint policy problem therefore stays infeasible even if the actor is told which mechanism applies. An unresolved source outcome for which no observation map is specified instead leaves the policy question unresolved; it is not another zero-policy result.

The checked finite method retains the full financial population for every mechanism and solves the shared-history constraints in two ways: enumeration of feasible instructions and an independent formulation of forbidden action combinations. It also compares the complete policy sets with explicit relabelings of the original policies. The checks distinguish a missing action row from an available safe action, and preserve the original temporal supports when equal numerical answers share storage.

These results explain how unresolved uncertainty can persist without requiring a separate explicit node for every time or transaction possibility. Symbolic temporal restrictions represent many possible times; shared numerical results avoid repeated calculation; common histories impose the information limits on decisions. Each operation preserves a different relation and needs its own justification. General observation languages, longer horizons, empirical actor access, saved partial frontiers and durable recovery require further work.\footnote{Complete-family binding, 26 family-change cases, 26 adversarial controls and saved independent replay are retained in \path{data/accounting-master/temporal-family-binding-20261027/}. The 21 cross-mechanism policy cases, independent full-set comparisons and corrected late-revelation expectation are retained in \path{data/accounting-master/cross-mechanism-policy-20261027/}. All message mappings and visibility regimes are explicit synthetic assumptions; neither reference introduces probabilities, actual actor knowledge or accounting permission.}


\section{A visible label need not identify the observation mechanism}
\label{sim:bound-observation-map}

In the preceding example, a first message $H$ requires pay under one interpretation and draw under another. An additional label can resolve the conflict if it tells the actor which interpretation applies. Suppose, however, that both interpretations attach the same label, \emph{reviewed}, to that message. The actor sees a label but still cannot distinguish the two cases. Declaring a label visible and declaring a mechanism identified are different assumptions.

The distinction changes the feasible policy set. With distinct labels supplied before the first decision, the two four-policy problems admit 16 joint policies. With the same label supplied in both, the conflicting instructions remain and no common policy is feasible. A simulator that automatically inserts its internal mechanism identifier into the actor's message would report the first answer while pretending to solve the second problem.

The selected observation declaration therefore records the actual token presented to the actor. It permits distinct mechanisms to emit the same token. Source identities and temporal proof labels remain attached to the analyst's explanation, but do not enter the decision history unless the declaration supplies a separate observation with that meaning. The numerical comparison keeps these cases apart even though both could informally be described as a labeled source.

\subsection{A message map must use information available at that date}

Hiding internal identifiers is insufficient if the message itself encodes an unavailable fact. A map that takes the entire financial history as its input could send one message when a hidden borrowing right exists and another when it does not. The resulting policy would exploit that right even if the source observations never revealed it. Giving the message an innocuous name would not remove the additional information.

The finite reference avoids that ambiguity by restricting what each map can read. At the first date, the coarse view supplies the market signal; the refined view supplies that signal and the current coupon-trigger observation. At the second date, the map receives the already specified trigger and cash-coverage observations. The actor retains the first message and its own earlier instruction. Hidden contractual rights, complete world identifiers and later observations cannot be introduced as new first-date map inputs.

For each selected source outcome, the declaration must map every supported observation key to a message. A missing key is an unfinished declaration, not permission to remove the corresponding financial history. Two entries for the same key are also invalid. These restrictions establish a closed finite interpretation of the existing observations. They do not establish that those observations describe a particular company's actual information practices, nor do they define a general language for arbitrary messages.

This ordering is useful when designing richer models. First specify which quantities can be observed and when; then specify how the observation is represented to the actor; finally constrain decisions to agree on indistinguishable histories. Reversing the order can make a desirable policy feasible by designing messages around the hidden state that policy would need to know.

\subsection{Source, actor and time remain separate premises}

A claimed visible tag carries an actor, the selected source revision, its content identity, a token and an observation time. The first decision must agree exactly with the temporal family's cutoff; the second occurs strictly later. A tag claimed at the first decision must have been available by that instant. Equivalent fixed-offset spellings describe the same instant, whereas a tag received one ten-millionth of a second afterward is late under the selected exact-time convention.

Unknown tag availability leaves the declaration unresolved. So does a supplied observation whose chronology cannot be established from the retained authorship bounds. The current reference treats this conservatively: it does not discard inconvenient clock possibilities or silently refine the source family with a new tag observation. A later method may condition or refine that family under explicitly admitted additional evidence, but it must preserve the distinction between the original question and the revised one.

These unresolved cases differ from a complete observation model with no feasible policy. In the latter, the available histories are specified and every possible instruction fails somewhere in the retained financial family. In the former, the model has not established which histories the actor may distinguish. Returning the same empty policy list for both would hide a missing information premise behind an apparently decisive financial conclusion.

A content identity detects alteration relative to expected bytes; it does not show that the actor received them. Likewise, the synthetic tag records used in the comparison test the consequences of supplied assumptions. They are not observations of a real company's managers or investors. That distinction remains necessary even when every mapping and timestamp check passes.

\subsection{Reuse the calculation only when the decision problem is unchanged}

A corrected source can preserve the same financial meaning and messages. The numerical policy calculation can then be reused while its source explanation changes. A different actor can also have the same supplied observation map; the actor binding must change even if the resulting mathematical problem is identical. Neither case licenses removal of the old and new source supports from the retained temporal family.

Changing a message label is different. A policy is a rule for responding to the labels the actor receives. Relabeling $H$ as another token changes that rule's interface even if four policies remain feasible and a mathematician can construct a one-to-one relabeling between the sets. The selected reference recomputes this case. It does not claim to have proved and implemented a general translation between equivalent policy representations.

The reuse comparison consequently includes the complete financial and action population, the actual first and second messages, the meanings of the actions, and the numerical methods. Identical repeated constraints can share computation while keeping every source outcome and temporal proof support attached. Distinct problems with zero feasible policies retain distinct identities. A future correction that is unavailable at the earlier cutoff remains in the source record without replacing the earlier view.

The independent checks reconstruct the messages from the supplied declarations and original observations, compare full policy sets from two formulations, and reject missing or duplicated keys, hidden-state inputs, wrong actors or times, altered sources and forged prior results. They also check changed second-date labels without changing first-date messages. These examples make the dependency rule concrete: identical counts are not enough, but a changed source explanation need not force an unchanged numerical problem to be solved again.

The remaining computational difficulty is how to continue an unfinished temporal search while retaining its unexamined possibilities. The conservative full-domain remainder already prevents a partial result from masquerading as complete; a checked, resumable partition would additionally preserve useful completed work. Actual message provenance, richer observation languages and refinement by new tag evidence remain distinct extensions.\footnote{The finite declarations, complete mapping comparisons, 22 main cases and 30 adversarial controls are retained in \path{data/accounting-master/observation-map-binding-20261027/}. A separate saved audit rechecks all main cases and four additional source, exact-time and token cases. The reference retains full source supports under numerical sharing and makes no claim of actual actor knowledge or authenticated evidence.}


\section{Stopping a search without losing the possibilities it has not examined}
\label{sim:temporal-frontier}

A correction may have reached the manager before a payment decision, or afterward. An interval for the delivery time can preserve both possibilities without assigning either a probability. If the source record contains several corrections, public releases and private deliveries, the exact temporal calculation may exhaust its work allowance before checking every combination. The unfinished calculation still has a useful result, provided it keeps the completed conclusions together with an exact description of what remains unexamined.

The conservative earlier method returned the whole original clock domain whenever its search stopped. That answer preserved uncertainty but discarded useful completed proof work. A resumable calculation can retain more. It must save the partition it has already established, including every unfinished part. Saving only the branches that produced interesting financial views would change the question on the next run.

\subsection{Each split must cover the entire timing domain}

Consider one cell $D$ of admissible event times under the supplied chronology and access assumptions. For the next source revision, let $g_1,\ldots,g_m$ be its eligibility conditions. A condition may require that publication or actor access occurred no later than the decision, or that ingestion occurred no later than the system's reconstruction date. The revision is eligible only if all these conditions hold.

A complete split keeps one eligible cell and one cell for each possible first failed condition:
\begin{align}
 D_{\mathrm{yes}} &= D\cap\bigcap_{j=1}^{m}g_j,\\
 D_{\mathrm{no},k} &= D\cap\bigcap_{j<k}g_j\cap g_k^{\mathrm c},
 \qquad k=1,\ldots,m.
\end{align}
Here a condition also denotes the set of time assignments satisfying it, and $g_k^{\mathrm c}$ denotes its exact complement. Every assignment in $D$ either satisfies every condition or has a unique first failure. These alternatives therefore cover $D$ and are pairwise disjoint. Their union is the original cell, not a sample of it. Some children may be empty, and a contradiction proof can establish that fact later.

The word \emph{exact} matters at a decision boundary. If $g_k$ says that access occurred at or before time $d$, its complement says strictly after $d$. Replacing that complement by \emph{at or after} would put the boundary in two children. Replacing the original condition by a strict inequality would remove the boundary from the eligible child. The selected reference retains rational event times and strict inequalities, so the preceding submicrosecond and fixed-offset examples keep their original meanings.

The search resolves one pending cell at a time. It first checks whether the cell is empty. A nonempty cell with no remaining revision to classify yields a selected or unresolved source view. Otherwise it installs the complete family of children above. Only after all those children exist may the calculation pause. This logical rule prevents a work limit from turning a complete split into a partial list of favorable branches.

\subsection{Pending does not mean feasible, and known does not mean complete}

A saved node has one of four meanings: its domain remains unchecked; a contradiction proves it empty; a witness establishes a terminal source view; or a witness supports an expanded cell whose children form the complete partition. An unchecked child need not contain a possible time assignment. It may already be inconsistent with an earlier chronology restriction, but the calculation has not yet established that. Retaining it is conservative; deleting it without checking is unsupported.

In the selected example, resolving the root installs eight children. Seven are later proved empty. With a one-node work allowance, all eight remain explicitly pending. The result therefore neither asserts eight feasible worlds nor pretends that the single eventual view is already established. The eight cells are a representation of unresolved work under known restrictions.

Once some terminal views are known, they can be displayed together with their source supports. Agreement among them still does not establish that the whole family has that view. An unexamined cell could contain a conflicting correction, a retraction or a different access history. Completion requires that no pending cell remains, including cells that eventually prove empty.

The same distinction controls policy interpretation. A policy that fails in an already established view fails for any larger family containing that view. A policy that succeeds in every established view may still fail in an unexamined one. Thus an incomplete temporal search can sometimes establish a counterexample before it can establish universal feasibility. Connecting that asymmetry to the numerical policy calculation requires an explicit treatment of the pending remainder; a saved search alone does not supply it.

\subsection{Resume the original question and preserve its proofs}

The checkpoint binds the complete source request, actor, trust declarations, timing assumptions, method identities and execution allowance. On reading it, a separate checker reconstructs each branch's constraints from the original domain and its path. It does not accept a saved mutable constraint list as a substitute. It verifies all completed witnesses or contradictions, the exact complement structure, source outcomes, support paths and work counts.

The reference uses a fixed left-first depth-first order. This makes the resumed tree directly comparable with the existing uninterrupted calculation. It also allows the checker to reject a purported completed branch placed after an earlier unfinished one. The order is a comparison convention, not a claim that this schedule is the fastest way to find useful financial conclusions.

A step allowance specifies how much new work one call may perform. A separately bound total execution cap determines when further work must stop. Reaching that cap retains the pending tree and reports a capped calculation; it does not turn unresolved cells into completed results. The test with a cap of two preserves the same original question and performs no repeated solver work when resumed at that cap. Changing the cap or the source assumptions requires an explicitly different, checked request to the continuation interface.

Across the 47 selected temporal requests, step allowances of one, two, seven and 4,096 produce exactly the same final proof trees, outcomes and source supports as uninterrupted search. Earlier certificates remain unchanged at every pause. A separate reread checks 1,060 saved checkpoints and rejects altered partition cells, proofs, source premises and completion declarations. These comparisons establish the selected finite continuation property.\footnote{The implementation, independent structural checker, 188 complete comparisons, 43 main adversarial controls and five additional saved-record controls are retained in \path{data/accounting-master/temporal-frontier-20261027/}. The main comparison makes 10,146 assertions; the saved-record audit makes 3,191. The source locks and manifests preserve the exact inputs, methods and comparison scope.}

This representation avoids a separate explicit node for every real-valued event-time assignment. It does not remove exponential growth in the number of distinguishable revision combinations. Nor does installing all children before a logical pause establish crash-safe file storage, authenticated evidence or actual actor knowledge. Larger revision families, better schedules, justified approximate reduction and durable storage require their own arguments and comparisons. The useful advance here is narrower: unfinished computation retains both what has been proved and every domain that has yet to be checked.


\section{Fixing the policy question before the search finishes}
\label{sim:prospective-partial-policy}

A manager who may pay an obligation or draw on a facility needs a rule for acting on the information available at each decision. The preceding search can establish some possible information histories while leaving others unexplored. To compare policies during that search, their meaning must remain fixed. If each newly discovered history changes what a policy is allowed to observe, an apparent reduction in the feasible set may instead reflect a change in the question.

The first requirement is therefore a declaration of possible messages before temporal search. In the finite example, start with the original registry of financial assertions and observation schemes. For each assertion-valued observation scheme, declare the first- and second-decision messages for every raw observation obtainable from the original financial assertions. Cross those financial and observation assertions conservatively. Some combinations may later prove temporally impossible; including them at this stage avoids selecting a favourable source history in advance.

The declaration binds the manager, decision dates, source records and meanings of pay and draw. A second-decision key also retains the manager's own earlier action, as in the preceding policy example. Internal source-head identifiers remain evidence of support and do not become messages to the manager. A retraction supplies no replacement assertion payload, but it does not justify removing every message associated with the original assertion before its timing has been examined. This is a declared finite vocabulary, not a finding that a real manager received every listed message.

\subsection{A fixed domain makes partial restrictions comparable}

Let $U$ denote the declared finite set of decision keys, including the decision stage and the required own-action memory. A total rule assigns either pay or draw to every key in $U$. Let $\Pi_U$ be the set of these rules. Assignments to keys that are never reached still belong to a total rule; they have no effect on that rule's realised path. Consequently, several total rules can represent the same policy on the reachable histories.

At search stage $t$, let $H_t$ contain only the selected financial and information histories supported by completed temporal branches. For a history $h$, write $V_h(\pi)$ when following rule $\pi$ selects no refused action in that history. This condition uses the complete finite action rows for that selected history, including first-decision refusal and second-decision refusal after the first action. The rules surviving the known restrictions are
\begin{equation}
 C_t=\{\pi\in\Pi_U:V_h(\pi)\text{ for every }h\in H_t\}.
 \label{eq:prospective-policy-restriction}
\end{equation}
These are sets of decision rules, not monetary balances or probability distributions. The relation is a local definition under a fixed message declaration and the supplied action semantics. Unresolved source outcomes are retained separately; they cannot supply a selected history or establish its financial constraints.

Resuming the same checked search retains the previously completed histories. Thus $H_t\subseteq H_{t+1}$. A rule satisfying every condition for the larger set also satisfies every condition for its subset, so
\begin{equation}
 C_{t+1}\subseteq C_t.
 \label{eq:prospective-policy-monotone}
\end{equation}
This inclusion explains the asymmetric conclusions. Once known constraints eliminate every rule, adding more histories cannot restore a rule under the same declaration. If some rules survive, a pending branch may still eliminate them. In the coarse/refined correction example, the known constraints already eliminate all rules after four of the 25 temporal nodes have been resolved. The baseline instead retains four total rules and establishes conditional feasibility only when the search completes with selected source outcomes throughout.

\subsection{Three reasons that a surviving rule may still be unresolved}

Unfinished search, unresolved source meaning and an inadequate domain allowance are different limitations. A two-node search cap retains 15 pending branches in the baseline. Those branches have not all been shown to contain possible histories, but they must remain represented until checked. The surviving rules therefore remain candidates.

A search can also finish while a required assertion remains unavailable, retracted or in conflict with another head. Unknown access leaves four rules surviving the known constraints, yet one source outcome remains unresolved. Retraction and conflicting-head cases likewise do not become feasible merely because all temporal branches have been processed. An inconsistent clock system is different again: it admits no temporal assignment. The 512 unconstrained rule assignments in that example do not establish a feasible policy for a possible company history.

Finally, the conservative domain may exceed the stated computational allowance. The finite profile permits at most 12 binary decision keys, hence at most 4,096 total rules. Renaming only the low message in one correction produces 12 keys and eight surviving total rules. Renaming all correction messages requires 18 keys, so the reference retains the original request and returns an unresolved domain budget without truncating the keys or enumerating a preferred subset. The extra rule count in the smaller example is not evidence of better information or better economic performance.

The checked partial calculation binds each result to the original request, declaration, domain, methods and temporal frontier. It retains the supporting branches, unresolved source outcomes and pending temporal domains. Changing a message, source, actor or domain cannot reuse the old result as the answer to the new question. A future extension rule may justify some reuse, but that relation must itself be stated and checked.\footnote{The original-registry domain comparison is retained in \path{data/accounting-master/prospective-policy-domain-20261027/}; its 10 cases include explicit overflow and conflicting-head cases. The partial comparison in \path{data/accounting-master/prospective-partial-policy-20261027/} checks 12 cases at 110 frontiers, with independent full-rule enumeration and 31 refusal controls. The candidate takes no complete-family comparison as an input. Exact sets and endpoint projections, rather than the counts alone, support these finite conclusions.}

This construction joins a fixed decision problem to an unfinished search without treating uncertainty as missing arithmetic. Longer horizons, additional actions, changing declarations and broader observation languages still require their own comparisons.



\section{Changing a message without confusing the decision problem}
\label{sim:policy-message-refinement}

A manager may receive a differently worded notice carrying the same information, or a notice that distinguishes two situations previously described alike. These changes have different consequences for a policy. The preceding requirement to fix message meaning does not prohibit either change. It requires the simulator to state which relation connects the old decision problem to the new one, and to identify the new result accordingly.

A consistent one-to-one renaming preserves the original distinction between every pair of messages. The finite comparison checks that renaming across every raw observation in every original source head, including heads not selected at the current frontier. It also preserves actor, decision dates, action meanings, source identities and second-stage memory of the manager's first action. Each total rule then has exactly one renamed counterpart. The first- and second-action refusal conditions have the same truth values after the key positions are permuted, so the full surviving rule sets correspond exactly. This is the same finite decision problem under different spellings; it is not an expansion of what the manager knows.

\subsection{A refined notice allows the old distinction to be recovered}

Now suppose a correction distinguishes a low message that was previously shared with another source history. The new message identifies which of those situations occurred. A manager can ignore that distinction and follow the old rule, or potentially act differently in the two situations. To make this precise, each refined message must map to one old message, consistently across all original raw observations. Combining this message map with the unchanged decision stage and own-action memory gives a map from new policy keys to old keys.

Write $U_o$ and $U_n$ for the old and new finite key domains, and $e:U_n\to U_o$ for this erasure map. Require $e$ to reach every old key. An old rule $\pi_o$ assigns pay or draw to each old key. Its lifted rule $L_e\pi_o$ chooses $\pi_o(e(u))$ at a new key $u$: the manager acts as if the additional distinction had not been received. New keys that erase to the same old key form a \emph{fiber}. A rule constant on each fiber makes the same choice at all of those keys.

For the same checked temporal frontier, let $C_o$ and $C_n$ be the rules satisfying the old and new known constraints. The comparison preserves each underlying action row and its refusal meaning, and checks that erasing the new forbidden assignments gives exactly the old forbidden assignments. Under those conditions,
\begin{equation}
 \{L_e\pi_o:\pi_o\in C_o\}
 =\{\pi_n\in C_n:\pi_n(u)=\pi_n(v)
       \text{ whenever }e(u)=e(v)\}.
 \label{eq:policy-refinement-lift}
\end{equation}
To see the first direction, a lifted rule makes the old choice at each corresponding raw observation. A refused first action or refused pair of successive actions would therefore also have been refused under the old rule. Conversely, a new rule constant on fibers defines a unique old rule because every old key has a preimage. The same correspondence of refused assignments then shows that this old rule satisfies the old constraints. The relation concerns exact finite rule sets under supplied message semantics; it assigns neither probabilities nor monetary values to information.

\subsection{Additional rule assignments need a reached-action comparison}

In the correction-specific low-message example, the old declaration permits four total rules and the refined declaration permits eight. Lifting preserves all four old rules. The other four also produce distinct sets of reached action rows across the complete selected family. Here the new distinction therefore permits additional feasible behavior under the stipulated information design. In another example, collapsing the first messages leaves no feasible rule, while restoring the refined messages permits four. Neither comparison establishes that those messages are actually available to a real manager or that any policy maximises value.

The source-unknown examples show why counting total rules is insufficient. With conflicting source heads, a refinement increases the unconstrained assignments from 32 to 2,048, but there is no selected history on which to demonstrate a new action. Every rule has the same empty reached-row signature. Under inconsistent clocks, the corresponding counts are 32 and 512, yet no admissible time assignment exists. An empty signature is a set used for comparison, not a realised or feasible company behavior. With unknown access, additional choices can be exhibited on the selected completed histories while the overall source family remains unresolved.

The comparison therefore retains both the full rule sets and their reached-row signatures. Pending branches, unknown source outcomes and inconsistent clocks keep their original meanings. The refined result has its own declaration and domain identity, even when its feasible rules correspond exactly to an older result.\footnote{Exact renaming is checked in \path{data/accounting-master/policy-domain-transport-20261027/}: 22 frontier comparisons and 12 refusal controls. Erasure and lifting are checked in \path{data/accounting-master/policy-observation-refinement-20261027/}: 17 comparisons and 10 controls, with whole-set equality and explicit reached-action comparisons. Split or merged labels are outside the renaming relation; a false refinement direction, altered source or dropped pending branch is refused. Both profiles remain finite conditional design evidence.}

This gives two precise ways to compare changed declarations. Broader changes to sources, decision times, actions or the observation language need additional arguments before earlier policy work can be reused.



\section{A correction arriving between two decisions}
\label{sim:joint-stage-source-policy}

A manager can make the first payment before receiving a correction, then consult the corrected information before the next financing decision. The first action belongs to the earlier information history. Replacing both observations with the latest source view would rewrite that history; treating the two dates independently could invent a sequence that no delivery time permits. The simulator needs a joint sequence of source views before it can interpret the two-stage policy.

Consider the supplied correction from a refined observation scheme to a coarse one. Let the first decision be on 2 January and the second on 3 January. If the correction is received before both dates, both select the coarse scheme. If it arrives between them, the first selects refined and the second coarse. Arrival after both leaves refined at both dates. When its access-time interval spans both decisions, the joint family therefore contains
\begin{equation}
 \{(\mathrm{coarse},\mathrm{coarse}),\,
   (\mathrm{refined},\mathrm{coarse}),\,
   (\mathrm{refined},\mathrm{refined})\}.
 \label{eq:joint-source-view-pairs}
\end{equation}
Each pair lists the first and second selected source schemes. These are possible sequences under the supplied chronology, not probabilities. Each marginal date admits both schemes, so their Cartesian product would also include coarse followed by refined. That fourth pair is impossible here: the same correction cannot have arrived before the first decision yet remain unavailable at the later one. A later retraction or a different replacement would be another source event requiring its own representation.

\subsection{Keep the source clocks shared while changing the decision cutoff}

The joint calculation uses the same source records, actor, access clocks and chronology at both dates. Its system reconstruction cutoff is fixed and may lie after both decisions; that cutoff describes the evidence inventory used by the analyst, not information available to the manager at the time. Each stage separately tests authorship, publication and actor access against its own decision date and applies its effective cutoff. The first increment keeps one financial assertion fixed while observation schemes can change.

The exact temporal partition first classifies eligibility at the first date and then at the second, carrying the accumulated clock constraints throughout. A completed branch records both selected or unresolved source views. An independent checker reconstructs the constraints and both head selections. The resulting joint pairs agree with the authored before, between and after cases, including exact decision boundaries and a submicrosecond arrival immediately after the first decision. Their projections agree with independently checked marginals, while the joint counterexample above tests the dependence that marginal agreement alone cannot establish.

Late retraction and conflicting heads leave the first view selected and the second unresolved. Unknown access retains every timing sequence admitted by the supplied chronology. It does not assert when the manager actually received the correction. Inconsistent clocks admit no joint sequence. If the joint search exhausts its allowance, this first profile retains the entire original clock question as unresolved; it does not yet save a resumable joint frontier.\footnote{The exact shared-clock comparison in \path{data/accounting-master/two-decision-source-views-20261027/} checks 11 cases and 12 refusal controls, preserving full joint partitions, source support and separately checked marginal results. Its 208 assertions include the impossible Cartesian-product sequence and the conservative budget remainder.}

\subsection{Use the selected map at each stage}

A prospective policy domain must now cover more than same-head combinations. Before inspecting completed joint outcomes, the declaration crosses the original financial head with every assertion-valued first-stage observation head and every assertion-valued second-stage head. It constructs the first message from the first selected map and the second message from the second selected map. The second decision key retains the actual first message and the manager's own first action. Internal head identities remain support records, not messages visible to the manager.

In the inherited raw observation schema, coarse versus refined changes only the first observation. The second raw observation remains the trigger and cash-coverage pair. Selecting a different source head at the second date therefore does not, by itself, add second-stage information. An explicit change to a second-message token tests whether the correct stage map is used, while remaining a supplied interface assumption.

That changed-token case needs 15 keys in the conservative original domain, exceeding the 12-key allowance. It remains unresolved without dropping mixed-stage combinations. A separate diagnostic collapses the first messages and uses 7 keys. Its encoded second messages differ when the wrong first-stage head is substituted for the second-stage head, so the test detects that wiring error. The correctly wired diagnostic nevertheless has no feasible policy. Correct source-to-message binding and policy feasibility are separate conclusions.

With the original mappings, arrival between the decisions or after both leaves 16 total rules on the conservative 11-key domain. Some assignments concern keys not reached in the selected family; these are not 16 demonstrated reached behaviors. Arrival before both dates supplies the coarse first view and leaves no feasible rule. The interval containing all three sequences also has no common rule because it includes that coarse-first history. Unknown access preserves the same conditional family and its infeasibility, without turning unknown access into an observed fact.

A required retracted or conflicting source view leaves the overall policy question unresolved, even when candidate assignments satisfy every known numeric constraint. Inconsistent clocks and either temporal or policy-domain budget exhaustion supply no feasible-policy answer.\footnote{\path{data/accounting-master/joint-stage-policy-20261027/} retains 16 cases,204 assertions and 7 refusal controls. A separate raw-observation mapper checks complete encoded rows; an independent policy oracle checks full rule sets. Same-head cases reproduce the earlier mapping, and the stage-specific token diagnostic tests preserved first-action memory. The 15-key case and the separately collapsed 7-key case are both retained.}

These results bind a complete joint source family to two-stage decisions. Continuing an unfinished joint search, changing financial rules between dates and admitting actual actor information remain separate tasks.



\section{What a paused joint search can establish}
\label{sim:joint-source-frontier}

Suppose the search stops after checking some of the correction's possible arrival times. A financing rule may survive those cases and fail in an unexamined case. To interpret the partial answer, the modeller needs both the source sequences already checked and the clock assignments still awaiting examination. The shared-clock construction in the preceding section supplies the latter without choosing a separate history at each date.

Each unfinished branch retains the constraints accumulated on the original clock variables. At a branch concerning a particular source and decision date, the alternatives are that every eligibility condition holds or that a named condition is the first to fail. These alternatives cover the parent branch without overlap. The search creates all of them before pausing. It can therefore resume from the first pending branch while preserving every other pending branch, the first decision's eligibility choices and the second decision's remaining tests. A short allowance changes how much has been examined, not the question being asked.

The saved state distinguishes a per-call allowance from a total node cap. Reaching the total cap leaves a capped, incomplete result with its pending domains. A checker reconstructs each branch's constraints from its path, verifies the rational feasibility or contradiction evidence and recomputes both selected source views at completed leaves. Checking only the number of leaves would miss a lost or duplicated support region. In the bounded examples, resuming one node at a time, seven at a time or without intermediate pauses produces the same complete tree, joint outcomes and support paths.

\subsection{Surviving rules can shrink before the search is complete}

Keep the prospective message domain fixed before examining these outcomes. Let $\Pi$ be its finite set of total two-stage decision rules: each rule assigns an action at every declared first- or second-stage message key, including keys that a particular source sequence might never reach. After $k$ units of search work, let $K_k$ contain the distinct fully selected source pairs found so far. For a pair $h$, let $F(h)\subseteq\Pi$ contain the rules satisfying every financial row supplied by that pair. The candidates consistent with the known pairs are
\begin{equation}
 C_k=\Pi\cap\bigcap_{h\in K_k} F(h),
 \qquad K_k\subseteq K_{k+1}
 \quad\Longrightarrow\quad C_{k+1}\subseteq C_k.
 \label{eq:joint-partial-policy-monotonicity}
\end{equation}
This is a local set-theoretic consequence of adding restrictions in the same rule universe. Initially no source pair may have been selected, in which case the empty intersection imposes no restrictions and $C_0=\Pi$. That initial survival says nothing about financing feasibility. Each newly selected pair can remove a rule; none can restore a rule already rejected by a known pair. The index $k$ measures search work, not elapsed financial time, so this relation does not describe how a manager's information evolves.

If the known restrictions leave no rule, no completion that merely adds further source pairs can supply a rule satisfying all of them. Conversely, a nonempty $C_k$ remains a partial answer while any source region is pending. In the correction example, a completed coarse-first pair can already remove all candidates, whereas observing only refined-first pairs cannot establish that no coarse-first pair remains. A changed source declaration or message domain would alter the comparison and requires a separate justification for reusing old results.

The returned conclusion also preserves the reason a family is unresolved. A retracted or conflicting required source remains an unresolved source view; the overall result keeps that status even if the known restrictions separately rule out every candidate. A complete timing family with unknown actual access is instead evaluated conditionally on all its supplied possible arrival sequences. Inconsistent clocks supply no admissible joint sequence and are not treated as a feasible empty family. Exceeding the prospective key allowance leaves the policy population unassessed. Only a completed search with selected source content and surviving rules supports conditional feasibility under these declarations.

\subsection{Check the reason for each rejected rule}

A saved list of surviving rules is insufficient evidence by itself. For every excluded total assignment, the result names a known financial row that refuses it. A first-stage refusal must match the rule's first action at that row's message. A second-stage refusal must also match the second action, using the second key that retains the actual first message and action. The checker validates these witnesses and requires the surviving and rejected assignments together to cover the entire fixed rule universe. It also reconstructs the conclusion, unknown source pairs, pending domains and support paths from the checked source search.

The finite comparison checks every saved endpoint and selected partial results against an independently evaluated policy population. Altered refusal witnesses, premature feasible conclusions, missing pending evidence and changed identities are rejected. The result checker still uses the established source-to-row mapping and policy oracle; it is not an independent redevelopment of all accounting and observation semantics.\footnote{The complete evidence is retained in \path{data/accounting-master/joint-temporal-frontier-20261027/}. Attempt 02 checks 11 requests, 33 completed search sequences, 662 saved checkpoints and 848 policy prefixes in 16 cases. Attempt 03 checks the 16 saved policy endpoints and 47 selected prefixes with 33 alteration controls, including wrong refusal witnesses and false feasible conclusions. The earlier failed import is preserved separately. These counts describe coverage of the supplied finite cases, not a general error rate.}

This establishes logical continuation for the supplied two-date question. It does not test crash-safe file storage or establish actual information delivery. Transporting a result to changed observation declarations, extending the horizon and admitting company evidence require further work.

